% Part: lambda-calculus
% Chapter: church-rosser
% Section: beta-reduction

\documentclass[../../../include/open-logic-section]{subfiles}

\begin{document}

\olfileid{lam}{cr}{b}

\olsection{Reduksi-$\beta$}

\begin{lem}\ollabel{lem:one-par}
  Yen $M \bredone M'$, mula
  $M \bredpar M'$.
\end{lem}
\begin{proof} Yen $M \bredone M'$
    lumantar kontraksi ing pucuk,
    $M$ awujud $(\lambd[x][N])Q$
    lan $M'$ awujud
    $\Subst{N}{Q}{x}$ kanggo
    sawetara $x$, $N$, lan~$Q$.
    Amarga $N \bredpar N$ lan
    $Q \bredpar Q$ miturut
    \olref[pb]{thm:refl},
    kita langsung oleh
    $(\lambd[x][N])Q
    \bredpar \Subst{N}{Q}{x}$
    saka \olref[pb]{defn:bredpar},
    mligine \olref[pb]{defn:bredpar4}.
    Kontraksi ing njero
    abstraksi utawa aplikasi
    uga kacakup dening aturan
    kompatibilitas reduksi paralel.
\end{proof}

\begin{lem}\ollabel{lem:par-red}
  Yen $M \bredpar M'$, mula
  $M \bred M'$.
\end{lem}

\begin{proof} Gunakake induksi ing !!{derivation}
  saka $M \bredpar M'$.
  \begin{enumerate}
  \item Yen aturan pungkasan
    \olref[pb]{defn:bredpar1},
    $M$ lan $M'$ mung $x$,
    lan $x \bred x$.
  \item Yen aturan pungkasan
    \olref[pb]{defn:bredpar2},
    $M$ awujud
    $\lambd[x][N]$ lan
    $M'$ awujud
    $\lambd[x][N']$
    kanggo sawetara $x$,
    $N$, $N'$, kanthi
    $N \bredpar N'$.
    Miturut hipotesis induksi,
    $N \bred N'$.
    Mula $\lambd[x][N] \bred
    \lambd[x][N']$ kanthi
    urutan kontraksi
    $\bredone$ kang padha
    karo $N \bred N'$.
  \item Yen aturan pungkasan
    \olref[pb]{defn:bredpar3},
    $M$ awujud $PQ$
    lan $M'$ awujud
    $P'Q'$ kanggo
    sawetara $P$, $Q$,
    $P'$, $Q'$, kanthi
    $P \bredpar P'$ lan
    $Q \bredpar Q'$.
    Miturut hipotesis induksi,
    $P \bred P'$ lan
    $Q \bred Q'$.
    Mula $PQ \bred P'Q'$
    kanthi nglakoni
    reduksi $P \bred P'$
    banjur reduksi
    $Q \bred Q'$.
  \item Yen aturan pungkasan
    \olref[pb]{defn:bredpar4},
    $M$ awujud
    $(\lambd[x][N])Q$
    lan $M'$ awujud
    $\Subst{N'}{Q'}{x}$
    kanggo sawetara $x$,
    $N$, $N'$, $Q$, $Q'$,
    kanthi $N \bredpar N'$
    lan $Q \bredpar Q'$.
    Miturut hipotesis induksi,
    $Q \bred Q'$ lan
    $N \bred N'$.
    Mula $(\lambd[x][N])Q
    \bred \Subst{N'}{Q'}{x}$
    kanthi nglakoni
    $N \bred N'$, banjur
    $Q \bred Q'$, lan
    pungkasane ngontraksi
    $(\lambd[x][N'])Q'$
    dadi $\Subst{N'}{Q'}{x}$.
  \end{enumerate}
\end{proof}

\begin{lem}\ollabel{lem:str}
  $\bred$ yaiku relasi
  transitif paling cilik
  kang ngemot $\bredpar$.
\end{lem}

\begin{proof}
  Dadia $\xred$ relasi
  transitif paling cilik
  kang ngemot $\bredpar$.

  $\bred \subseteq \xred$:
  Upamane $M \bred M'$.
  Yen ana langkah kontraksi,
  iki ateges $M \ident M_1
  \bredone \dots \bredone
  M_k \ident M'$.
  Miturut \olref{lem:one-par},
  $M \ident M_1 \bredpar
  \dots \bredpar M_k
  \ident M'$.
  Amarga $\xred$ ngemot
  $\bredpar$ lan transitif,
  $M \xred M'$.
  Yen nol langkah, identitas
  wis kacakup dening
  refleksivitas reduksi paralel.

  $\xred \subseteq \bred$:
  Upamane $M \xred M'$,
  yaiku $M \ident M_1
  \bredpar \dots
  \bredpar M_k \ident M'$.
  Miturut \olref{lem:par-red},
  $M \ident M_1 \bred
  \dots \bred M_k \ident M'$.
  Amarga $\bred$ transitif,
  $M \bred M'$.
\end{proof}

\begin{thm}\ollabel{thm:cr}
  $\bred$ nduweni
  sipat Church--Rosser.
\end{thm}

\begin{proof}
  Langsung saka
  \olref[dap]{thm:str},
  \olref[pb]{thm:cr},
  lan \olref{lem:str}.
\end{proof}

\end{document}
